\chapter{Single Page Application}

Abbiamo imparato a sviluppare applicazioni web tradizionali, con CGI, PHP,
Node.js puro e con un framework. Nelle applicazioni che abbiamo costruito le
pagine sono generate dinamicamente sul server e inviate al browser per la
visualizzazione: una volta arrivate al browser, però, sono intrinsecamente
statiche.

\section{Applicazioni multi-pagina}

Nelle nostre applicazioni tradizionali, quando si cambia pagina (per esempio
facendo clic su un collegamento):

\begin{itemize}
  \item viene inviata una nuova richiesta al server;
  \item la nuova pagina viene restituita nella risposta e disegnata dal browser.
\end{itemize}

Questo approccio è noto come \term{multi-pagina}, e porta ad applicazioni web
multi-pagina. Navigare verso una pagina diversa provoca sempre un
\textbf{ricaricamento completo}.

Nel capitolo 7 avevamo studiato JavaScript nell'ambiente del browser:
manipolazione del DOM, recupero di dati, eventi. Finora, però, non abbiamo fatto
un uso particolare di quelle funzionalità.

\section{Single Page Application}

\dfn{Single Page Application}{
  Una \term{Single Page Application} (SPA) reagisce all'input dell'utente
  \textbf{aggiornando dinamicamente il contenuto della pagina corrente}, invece
  di caricare una nuova pagina dal server. L'intera applicazione viene caricata
  con la richiesta iniziale, e da quel momento la pagina non viene più
  ricaricata.}

La complessità della generazione delle viste si sposta dal server al codice
lato client: da \textit{thin client} a \textit{thick client}. Sono applicazioni
a forte intensità di JavaScript.

\begin{figure}[H]
  \centering
  \begin{tikzpicture}[font=\Lato\scriptsize,
      n/.style={rounded corners=3pt, line width=0.8pt, align=center,
                minimum width=2.5cm, minimum height=0.85cm},
      c/.style={n, draw=accent!60!black, fill=accent!10!white},
      s/.style={n, draw=primary!65!black, fill=primary!10!white},
      a/.style={-{Stealth[length=5pt,width=4pt]}, line width=0.8pt,
                draw=neutral!45!black}]

    % multipagina
    \node[font=\Lato\scriptsize\bfseries, text=neutral!25!black] at (2.6,2.1)
      {multi-pagina};
    \node[c] (b1) at (0,1.2) {Browser};
    \node[s] (s1) at (5.2,1.2) {Server};
    \draw[a] ([yshift=5pt]b1.east) -- node[wtlab, above]{richiesta}
             ([yshift=5pt]s1.west);
    \draw[a] ([yshift=-5pt]s1.west) -- node[wtlab, below]{pagina HTML completa}
             ([yshift=-5pt]b1.east);
    \node[font=\Lato\scriptsize\itshape, text=red-gradient-6!75!black]
      at (2.6,0.3) {ricaricamento completo a ogni navigazione};

    % SPA
    \node[font=\Lato\scriptsize\bfseries, text=neutral!25!black] at (2.6,-1.3)
      {single page};
    \node[c] (b2) at (0,-2.2) {Browser};
    \node[s] (s2) at (5.2,-2.2) {API REST};
    \draw[a] ([yshift=5pt]b2.east) -- node[wtlab, above]{richiesta di dati}
             ([yshift=5pt]s2.west);
    \draw[a] ([yshift=-5pt]s2.west) -- node[wtlab, below]{solo JSON}
             ([yshift=-5pt]b2.east);
    \node[font=\Lato\scriptsize\itshape, text=green-gradient-6!60!black]
      at (2.6,-3.1) {il JavaScript lato client ridisegna solo ciò che cambia};
  \end{tikzpicture}
  \caption{Multi-pagina contro single page.}
  \label{fig:spa-vs-mpa}
\end{figure}

Navigare verso una \guillemotleft pagina\guillemotright\ diversa non provoca un ricaricamento
completo: i dati necessari vengono recuperati da un backend (per esempio una API
REST) e il JavaScript lato client costruisce dinamicamente la nuova pagina.

\section{Routing lato client}

Le diverse \guillemotleft pagine\guillemotright\ di una SPA si chiamano tipicamente \term{viste}. La
navigazione e l'aggiornamento delle viste sono gestiti \textbf{interamente} lato
client, senza ricaricare la pagina dal server: è il \term{routing lato client}.

Comporta tre cose:

\begin{itemize}
  \item aggiornare la URL nella barra degli indirizzi;
  \item caricare e disegnare il contenuto (la vista) appropriato;
  \item interagire con la \term{History API} del browser, perché i pulsanti
        \guillemotleft indietro\guillemotright\ e \guillemotleft avanti\guillemotright\ continuino a funzionare come l'utente si
        aspetta.
\end{itemize}

I componenti chiave sono quindi un \textbf{router} (che mappa le URL sulle
viste), un meccanismo di \textbf{rendering delle viste} e una
\textbf{gestione della cronologia}.

\section{Una SPA da zero}

\subsection{L'unico documento HTML}

C'è un solo documento HTML, e non contiene alcun contenuto: soltanto un
\tg{div} vuoto. È \code{main.js} a fare tutto il lavoro.

\begin{lstlisting}[style=html, name=index.html]
<!DOCTYPE html>
<html lang="it">
  <head>
    <title>Single Page Application</title>
    <link rel="stylesheet" href="src/css/style.scss"/>
  </head>
  <body>
    <div id="app"></div>
    <script type="module" src="./src/js/main.js"></script>
  </body>
</html>
\end{lstlisting}

\subsection{La struttura di base}

Lo script comincia creando tre elementi dentro il contenitore \code{app}: una
barra di navigazione, un contenitore per il contenuto principale e un piè di
pagina. La barra e il piè di pagina non cambieranno passando da una vista
all'altra; il contenuto sì.

\begin{lstlisting}[style=js, name=main.js]
let nav = document.createElement("nav");
nav.id = "nav";

let content = document.createElement("main");
content.id = "content";

let footer = document.createElement("footer");
footer.id = "footer";

app.append(nav, content, footer);
\end{lstlisting}

\subsection{Le rotte}

Ogni rotta associa una vista a una data URL, ed è indicizzata dal proprio
percorso. A ciascuna sono associati un \code{title} (che diventerà il titolo
della pagina) e una funzione \code{render} che restituisce una stringa HTML.

\begin{lstlisting}[style=js, name=main.js]
import about from "./views/about.js";   // altre importazioni omesse

const routes = {
  "/":         { title: "Home",             render: home },
  "/about":    { title: "About",            render: about },
  "/webtech":  { title: "Web Technologies", render: contact },
};
\end{lstlisting}

\begin{lstlisting}[style=js, name=views/about.js]
export default () => /*html*/`
  <h1>About</h1>
  <p>
    Questa e' una single page app. Navigare verso pagine diverse non
    provoca un ricaricamento completo!
  </p>
`;
\end{lstlisting}

La generazione della vista è ridotta all'osso: una funzione che restituisce una
stringa HTML. Può diventare molto più complessa --- tanto per cominciare, si
potrebbe usare un motore di template.

\subsection{Il router}

\begin{lstlisting}[style=js, name=main.js]
function router() {
  let view = routes[location.pathname];   // rotta per il percorso corrente

  if (view) {
    document.title = view.title;          // aggiorna il titolo della pagina
    content.innerHTML = view.render();    // aggiorna il contenuto
  } else {
    // nessuna rotta corrisponde: sostituisce la voce di cronologia e va a "/"
    history.replaceState("", "", "/");
    router();
  }
}
\end{lstlisting}

\subsection{I collegamenti di navigazione}

\begin{lstlisting}[style=js, name=main.js]
for (let route in routes) {
  let link = document.createElement("a");
  link.href = route;
  link.innerText = routes[route].title;
  link.setAttribute("data-link", "");   // attributo personalizzato
  nav.append(link);
}

footer.innerHTML = renderFooter();
\end{lstlisting}

In una SPA possono coesistere collegamenti che devono comportarsi normalmente
(per esempio verso siti esterni) e collegamenti che devono soltanto provocare un
cambio di vista. Nell'esempio i secondi sono identificati dall'attributo
\code{data-link}.

\begin{lstlisting}[style=js, name=main.js]
window.addEventListener("click", e => {
  if (e.target.matches("[data-link]")) {
    e.preventDefault();                        // niente ricaricamento
    history.pushState("", "", e.target.href);  // nuova voce di cronologia
    router();                                  // aggiorna la vista
  }
});
\end{lstlisting}

Si riconosce qui la \textbf{delegazione degli eventi} studiata nel capitolo 7:
un solo gestore registrato su \code{window} serve tutti i collegamenti, compresi
quelli creati dinamicamente in seguito.

\subsection{Sincronizzarsi con la cronologia}

Resta da invocare \code{router()} per allineare la vista al percorso corrente
in due situazioni: quando l'utente preme il pulsante \guillemotleft indietro\guillemotright\ (evento
\code{popstate}) e quando la pagina viene ricaricata per intero (evento
\code{DOMContentLoaded}).

\begin{lstlisting}[style=js, name=main.js]
window.addEventListener("popstate", router);
window.addEventListener("DOMContentLoaded", router);
\end{lstlisting}

\warning{Il server deve collaborare}{
  Se l'utente ricarica la pagina mentre si trova su \code{/about}, il browser
  invia al server una richiesta per \code{/about} --- che sul server non esiste
  come file. Il server di una SPA va quindi configurato per restituire sempre
  \code{index.html} per qualunque percorso, lasciando che sia il router lato
  client a decidere quale vista mostrare.}

\section{I componenti}

Man mano che la complessità cresce, gestire intere viste come monoliti diventa
rapidamente impraticabile: il codice si fa difficile da mantenere ed è
complicato riusare pezzi di funzionalità in viste diverse.

La strada è \textbf{scomporre l'interfaccia in componenti modulari}. Ogni
componente incapsula una funzionalità specifica ed è l'\textbf{unico
responsabile} del disegno della propria interfaccia e della gestione del proprio
stato interno; le viste che lo usano non devono preoccuparsi di questi
dettagli.

\begin{center}
\renewcommand{\arraystretch}{1.4}
\begin{tabular}{@{}l@{\hspace{1.2cm}}>{\raggedright\arraybackslash}p{8.6cm}@{}}
  \textbf{Caratteristica} & \textbf{Significato} \\
  \term{Modularità}    & devono essere autocontenuti e incapsulare una
                         funzionalità specifica \\
  \term{Incapsulamento}& devono incapsulare la propria logica, il proprio stato,
                         il proprio layout e i propri stili \\
  \term{Riusabilità}   & devono essere progettati per essere riusati in parti
                         diverse dell'applicazione \\
\end{tabular}
\end{center}

I componenti hanno un ruolo centrale in tutti i framework di front-end.

\begin{esempio}{Un componente con le Web Components API}
\begin{lstlisting}[style=js, name=components/counter.js]
let count = 0;

class Counter extends HTMLElement {
  constructor() {
    super();
    this.innerHTML = `<button>Clicks : ${count}</button>`;

    let btn = this.querySelector("button");
    btn.onclick = () => {
      btn.innerHTML = "Clicks : " + ++count;
    };
  }
}

customElements.define("click-counter", Counter);
\end{lstlisting}

Stiamo definendo un nuovo \code{HTMLElement} chiamato \code{<click-counter>},
associato alla classe \code{Counter}. Il browser userà quella classe ogni volta
che incontrerà un elemento \code{<click-counter>} nel DOM.

\begin{lstlisting}[style=js, name=views/home.js]
import "../components/counter.js";

export default () => /*html*/`
  <h1>Home</h1>
  <p>Semplice contatore di clic</p>
  <click-counter></click-counter>
`;
\end{lstlisting}

Dal punto di vista della vista, il componente è indistinguibile da un tag HTML
qualunque: è precisamente l'incapsulamento che cercavamo.
\end{esempio}

\section{Perché una SPA}

Le SPA offrono alcuni vantaggi rispetto alle applicazioni multi-pagina
tradizionali:

\begin{itemize}
  \item un'esperienza utente \textbf{più dinamica}, simile a quella di
        un'applicazione desktop;
  \item il caricamento iniziale è più lento (bisogna scaricare l'intera
        applicazione), ma quelli successivi sono \textbf{molto} più rapidi;
  \item \textbf{carico ridotto sul server}: dal server si recuperano solo i dati
        necessari e solo quando servono, senza trasferire ogni volta interi
        documenti HTML, con risparmio di banda.
\end{itemize}

\begin{infobox}{{Una nota su AJAX}}
  Prima che la \code{fetch} API fosse introdotta con ES6, le richieste di rete
  si effettuavano con \term{AJAX} (\textit{Asynchronous JavaScript And XML}),
  implementato tramite l'interfaccia \code{XMLHttpRequest}. Oggi
  \code{XMLHttpRequest} è largamente sostituito da \code{fetch}, ma è ancora
  piuttosto comune e lo si può incontrare in alcuni pacchetti.

\begin{lstlisting}[style=js]
let xhttp = new XMLHttpRequest();
xhttp.onreadystatechange = function () {
  if (this.readyState == 4 && this.status == 200) {
    document.getElementById("demo").innerHTML = xhttp.responseText;
  }
};
xhttp.open("GET", "/url/to/fetch", true);
xhttp.send();
\end{lstlisting}
\end{infobox}

\section{I limiti del fare tutto a mano}

Realizzare la lista di cose da fare come SPA in JavaScript puro è stato
istruttivo, ma nonostante la semplicità dell'applicazione sono emerse alcune
difficoltà.

\begin{center}
\renewcommand{\arraystretch}{1.45}
\begin{tabular}{@{}l@{\hspace{1.0cm}}>{\raggedright\arraybackslash}p{8.6cm}@{}}
  \textbf{Problema} & \textbf{Perché è difficile} \\
  \textbf{Gestione dello stato}
    & bisogna tenere stato e interfaccia sempre allineati, e ricordarsi di
      aggiornare le viste dopo ogni cambiamento di stato \\
  \textbf{Manipolazione del DOM}
    & manipolare il DOM con l'API \code{document} richiede molto codice
      ripetitivo \\
  \textbf{Routing lato client}
    & il nostro router era elementare: e se servissero rotte parametriche,
      reindirizzamenti, rotte protette? \\
  \textbf{Rendering efficiente}
    & ridisegniamo l'intera lista a ogni modifica; per essere efficienti
      bisognerebbe aggiornare solo ciò che è cambiato \\
\end{tabular}
\end{center}

I \term{framework di front-end} come \term{Angular}, \term{React},
\term{Vue.js}, \term{Ember} e \term{Svelte} forniscono soluzioni già pronte per
queste difficoltà, e tipicamente includono anche supporto strumentale:
generazione di codice ripetitivo, automazione di build e rilascio.

Ciascun framework fornisce inoltre un modo in qualche misura
\guillemotleft standardizzato\guillemotright\ di costruire front-end: chi è abituato a lavorare con
un dato framework avrà meno difficoltà a inserirsi in un altro progetto che usa
lo stesso framework. Nel prossimo capitolo conosceremo Angular.

\section{Riferimenti}

\begin{itemize}
  \item \textbf{SPA (Single-page application)} --- MDN web docs. \\
        \urlx{https://developer.mozilla.org/en-US/docs/Glossary/SPA}
  \item \textbf{Single page apps in depth}, Mikito Takada. \\
        \urlx{http://singlepageappbook.com/} \\
        Parte rilevante: \textit{Modern single page apps --- an overview}.
  \item \textbf{Understanding client-side JavaScript frameworks} --- MDN web
        docs. \\
        \urlx{https://developer.mozilla.org/en-US/docs/Learn/Tools_and_testing/Client-side_JavaScript_frameworks} \\
        Parti rilevanti: \textit{Introduction to client-side frameworks} e
        \textit{Framework main features}.
\end{itemize}
